% ======================================================================
%  Section 4.7.3 --- Weight maps.  Corrected.
%
%  The argument of this section is right and it is the section that
%  carries Variant 3 now that 4.7.2 reports a tie on accuracy. But your
%  own Table 4.3 contradicts one of its central claims, and correcting
%  that turns out to explain the nominal penalty of 4.7.1. Nine items.
%
%  --------------------------------------------------------------------
%  A. RESOLVED. These maps are at iteration 600 of the occluded run, and
%     4.7.2 is now reported at iteration 600 as well, so the two sections
%     describe the same instant and nothing needs reconciling. This item
%     stands only as a record that it was checked.
%
%  --------------------------------------------------------------------
%  B. "THE WEIGHTING HAS BECOME, TO WITHIN A THIRD OF A PERCENT, THE
%     IDENTITY." Table 4.3 says otherwise, on the same line:
%         clean, iteration 600 : mean weight 0.789, and only 14.1% of
%         measurements at w >= 0.99 -- FEWER than the 24.7% at iteration 1.
%     A weighting whose mean is 0.789 and which leaves six measurements in
%     seven below unit weight is not the identity and is not close to it.
%     What falls from 3.5% to 0.3% is the OUTRIGHT REJECTION, which is a
%     different statement and the one the data support.
%
%     The reason it cannot become the identity is worth stating, because
%     it is structural rather than accidental: the scale is estimated from
%     the residuals, so it contracts as they do. The ratio e/s does not
%     shrink as the pose converges, and the Tukey function keeps assigning
%     intermediate weights however small the residuals become. An
%     M-estimator with an adaptive scale never switches itself off.
%
%     AND THIS EXPLAINS 4.7.1. The nominal table shows Variants 2 and 3
%     LESS accurate than the baseline on a clean scene, and 4.7.1 attributes
%     it to valid measurements being rejected. Now you can say why the
%     effect persists to convergence instead of vanishing with the pose
%     error: the weighting is still active at iteration 600, because it is
%     always still active. Three sections lock together once this
%     paragraph is corrected, so it is worth the rewrite.
%
%  --------------------------------------------------------------------
%  C. 7.4 PER CENT IS THE WRONG COVERAGE FOR ITERATION 600.
%     "9.5% is the 7.4% it covers together with some two percent of edge
%     response." The occluder covers 7.4% at the INITIAL pose and 5.80% at
%     the desired pose, which is where the camera is at iteration 600.
%     4.7.2 and 4.6.3 both state both figures. So the arithmetic reads
%     9.5% against 5.80%, a ratio of 1.6, and the excess is not two per
%     cent but three and a half. Rewritten, and see (D).
%
%  --------------------------------------------------------------------
%  D. THE ARGUMENT AT ITERATION 1 CAN BE MADE MUCH STRONGER.
%     You write that 4.2% rejected is less than the 7.4% covered, so the
%     occluder is not yet being rejected. True but weak: it allows that
%     four-sevenths of the occluder is rejected. The clean run gives you
%     the control. It rejects 3.5% with no occluder present at all, so the
%     occluder accounts for at most 4.2 - 3.5 = 0.7 per cent -- under a
%     TENTH of the region it covers. That is the rigorous form of "not yet
%     being rejected at all", and it costs one sentence.
%
%  --------------------------------------------------------------------
%  E. WHOSE WEIGHTS, AND HOW MANY? Two problems in one caption.
%     (i) "the Tukey weights of Variant 3". Variant 2 IS the Tukey
%         variant. Variant 3's weights are w^H of (4.13) -- the Tukey
%         function applied to structure-normalised residuals. Calling them
%         Tukey weights invites the reader to think the figure shows
%         Variant 2.
%     (ii) "over the 66,000 measurements of the interior region". There
%         are 264,000 measurements, four per position, and 4.4.3 defines
%         w_j one per MEASUREMENT. The four at a position share the same
%         Gbar but have different residuals e_j, so they have four
%         DIFFERENT weights. A map with one value per position is
%         therefore showing one scale, or an aggregate of the four, and
%         the caption must say which.
%     This also collides with my correction to 4.6.2, which states that
%     these percentages are fractions of P. If the statistics are over
%     66,000 then that sentence in 4.6.2 is wrong and I need to fix it --
%     tell me which the code computes and I will make them agree.
%
%  --------------------------------------------------------------------
%  F. THE MID-GREY BORDER IS INDISTINGUISHABLE FROM A REAL WEIGHT.
%     The caption assigns mid-grey to the filter border "so as not to be
%     read as rejection". But the mean weight is 0.768 to 0.789, so much
%     of the image IS mid-grey, legitimately. A reader cannot tell the
%     border from a measurement of weight one half. Crop the border out,
%     or mark it in a colour no weight can take.
%     Related: "white denotes retained, black denotes rejected" describes
%     a binary map. This one is continuous -- that is the whole content of
%     Table 4.3 -- and the caption should say the grey level is the weight.
%
%  --------------------------------------------------------------------
%  G. THE SELF-REFERENCE OF 4.7.1 AGAIN, TWICE. Both uses of
%     \ref{sec:ablation-results} mean 4.7.2 and should be sec:occlusion;
%     as written they point at the parent section 4.7.
%
%  H. "DISC" against the "elliptical region" of 4.6.3 and 4.7.2. Four
%     occurrences. Pick one word for the whole chapter.
%
%  I. THE DILATION IS MEASURED AND IS WIDER THAN THE FILTERS EXPLAIN.
%     Rejected 9.5% against 5.80% occluded is a band of roughly 8 pixels
%     around the occluder. The 9x9 filters have a half-width of 4, and a
%     filter centred more than 4 pixels from the boundary cannot see the
%     occluder at all, so support alone accounts for at most half the band
%     I measure. I have written the sentence to report the measurement and
%     the expectation without inventing a mechanism to bridge them, and
%     flagged it below as something to check by measuring the band
%     directly off the map. Do not let me assert an explanation here that
%     the measurement does not support -- it is exactly the kind of claim
%     we removed from 4.3.2 and 4.5.
% ======================================================================

\subsection{Weight maps}
\label{sec:weight-maps}

The pose errors of Sections~\ref{sec:nominal} and~\ref{sec:occlusion}
establish what the weighting achieves but not how. Since the weights are
defined one per measurement, they can be displayed at their image positions,
and doing so for both scene conditions at the beginning and the end of the
servo makes the mechanism directly observable.
Figure~\ref{fig:wmaps} collects the four maps: the grey level of each position
is the weight carried there, white denoting a measurement retained at
$w_{j}=1$ and black one rejected at $w_{j}=0$, and the ten-pixel border that
the $9\times9$ filters cannot reach carries no measurement and is excluded.
The accompanying statistics are collected in Table~\ref{tab:weight-stats}.
%%% (F) The map is continuous, not binary, and saying so costs nothing.
%%% The border is now described as excluded rather than as mid-grey; see
%%% the note at the foot on why mid-grey cannot be used.
%%% (G) sec:ablation-results -> sec:occlusion.

\paragraph{At the start of the servo the estimator cannot tell the two scenes
apart:}
At iteration~1 the clean and the occluded maps are almost indistinguishable:
the mean weight is $0.768$ against $0.751$, and the fraction of measurements
rejected outright is $3.5\%$ against $4.2\%$. In neither case does the
occluder stand out. The clean run provides the control that makes this
precise: it rejects $3.5\%$ with nothing present to reject, so of the $4.2\%$
rejected in the occluded run at most $0.7\%$ can be attributed to the
occluder --- under a tenth of the $7.4\%$ of the image it covers at this pose.
The occluder is, at the first iteration, very largely retained.
%%% (D) The comparison against the clean run replaces the comparison
%%% against 7.4%, and turns "less than" into "under a tenth".
The reason is the robust scale. With the initial misalignment of some
$41$~pixels the residual is large over the whole image, so the median
absolute deviation~\eqref{eq:mad} is large in proportion, and no group of
measurements is extreme relative to it. What rejection does occur is diffuse
and spread across the frame, and since the clean scene contains nothing to
reject, all $13.6\%$ of its strongly down-weighted measurements are valid
ones. This is precisely the difficulty anticipated in
Section~\ref{sec:tukey-limitation}, and the nominal map at iteration~1 is its
direct measurement.

\paragraph{By convergence the two have separated completely:}
At iteration~600 the clean scene rejects $0.3\%$ of measurements outright,
against $3.5\%$ at the first iteration, and holds only $3.2\%$ below a quarter
weight, against $13.6\%$: the map is white but for a faint tracery along the
strongest edges, where a residual misalignment of a fraction of a pixel still
registers. The rejection of valid measurements has all but ceased.

It has not, however, given way to the identity, and the distinction matters.
The mean weight at convergence is $0.789$, and only $14.1\%$ of measurements
carry a weight above $0.99$ --- fewer, not more, than the $24.7\%$ at the
first iteration. What has changed is the shape of the distribution rather
than its centre: the extremes at both ends have been drawn in, leaving the
bulk of the measurements at intermediate weights. This is structural and not
accidental. The scale is estimated from the residuals themselves, so it
contracts as they do, and the ratio $\tilde{e}_{j}/\hat{s}_{\tilde{e}}$ on
which the Tukey function acts retains a comparable spread however small the
residuals become. An M-estimator with an adaptive scale does not switch itself
off when the data cease to contain outliers; it continues to discriminate
among whatever residuals remain. An independent trace of the same fact is
reported in Section~\ref{sec:occlusion}, where Variant~3 retains a commanded
velocity thirty-four times that of Variant~2 after its pose has ceased to
improve: the weights are still changing at iteration~$600$, which they could
not be had the weighting settled to a constant matrix.
%%% (B) This replaces "the weighting has become, to within a third of a
%%% percent, the identity", which Table 4.3 contradicts on the same line.
%%% The cross-reference to the terminal jitter is new: it is a second,
%%% independent observation of the same thing, from a quantity logged for
%%% a different purpose, and the two together make the point unarguable.

The observation accounts for a result of Section~\ref{sec:nominal} that was
there reported without explanation. On the clean scene the weighted variants
position less accurately than the baseline, and the penalty does not vanish as
the pose error falls. It cannot vanish: the weighting is still active at
convergence because it is always still active, so a weighted variant continues
to discard information that the unweighted scheme retains, right down to the
pose at which the run stalls. The cost of robustness on a clean scene is
therefore permanent rather than transient, and it is the price paid for the
two orders of magnitude recovered under occlusion.
%%% This paragraph did not exist and is the reason to make the correction
%%% in (B) rather than merely soften the sentence. It closes the loop
%%% between 4.3.3, 4.7.1 and this section.

The occluded scene at the same iteration rejects $9.5\%$. The occluder appears
as a solid, sharply bounded black region, and the rejection is no longer
diffuse: it has concentrated onto the measurements that are corrupted. The
rejected area exceeds the $5.80\%$ that the occluder covers at this pose by a
factor of $1.6$, the excess forming a band around it some eight pixels wide.
Part of this is to be expected: a filter whose $9\times9$ support straddles
the boundary responds to the occluder although its centre lies outside it, so
a band of up to the four-pixel half-width of the support must be rejected with
it. The measured band is wider than that argument accounts for, and the
remainder is not explained here.
%%% (C) 7.4% -> 5.80%, since at iteration 600 the camera is at the desired
%%% pose. (I) The dilation is reported as measured, with the part that is
%%% explained separated from the part that is not. See the note at the
%%% foot for the measurement that would close it.

\paragraph{The discrimination is a product of convergence, not a precondition
for it:}
Read together, the four panels describe a mechanism that sharpens as the servo
proceeds. While the pose error is large the misalignment residual dominates,
the robust scale is large, and the estimator discriminates poorly in both
scenes alike. As the pose error falls the misalignment residual collapses, the
scale falls with it, and the occluder --- whose residual does not fall,
because no pose can cancel it --- becomes the only large residual in the image
and is rejected decisively. Robust weighting and pose convergence therefore
reinforce one another, rather than the weighting simply enabling the
convergence, and the effect is visible in the figure as the divergence of the
right-hand column from the left.

Two consequences follow for how the preceding sections are read. First, the
cost identified in Section~\ref{sec:tukey-limitation} has two parts of
different character. The over-rejection of valid measurements early in the
servo is transient and largely disappears, falling from $13.6\%$ to $3.2\%$ of
the clean scene between the first iteration and convergence; the residual
down-weighting, by contrast, persists, for the reason given above. Second, the
failure of Variant~1 under occlusion reported in Section~\ref{sec:occlusion}
is not attributable to any single catastrophic iteration but to the
accumulation of a bias that the weighted variants remove progressively, and
which they remove most effectively precisely when the remaining pose error is
smallest.
%%% The first consequence previously said the cost "is real but transient".
%%% After (B) that is only half true, and the half that is not true is the
%%% half that explains 4.7.1. Split into its two parts.
%%% (G) second \ref{sec:ablation-results} -> sec:occlusion.

\begin{table}[H]
    \centering
    \caption{Distribution of the weights $w^{\mathrm{H}}_{j}$ of Variant~3
    over the interior region, for the clean and occluded scenes at the first
    and last iteration. The occluder covers $7.4\%$ of the measurements at the
    initial pose and $5.80\%$ at the desired pose.}
    \label{tab:weight-stats}
    \small
    \begin{tabular}{@{}llrrrr@{}}
        \toprule
        Scene & Iteration & $\bar{w}$ & $w_j\le0.01$ & $w_j<0.25$ & $w_j\ge0.99$ \\
        \midrule
        clean    & $1$   & $0.768$ & $3.5\%$ & $13.6\%$ & $24.7\%$ \\
        clean    & $600$ & $0.789$ & $0.3\%$ & $\phantom{0}3.2\%$ & $14.1\%$ \\
        occluded & $1$   & $0.751$ & $4.2\%$ & $16.0\%$ & $26.8\%$ \\
        occluded & $600$ & $0.719$ & $9.5\%$ & $14.2\%$ & $17.8\%$ \\
        \bottomrule
    \end{tabular}
\end{table}
%%% (E)(i) "Tukey weights of Variant 3" -> the weights w^H, since Variant 2
%%% is the Tukey variant and the old wording points at the wrong figure.
%%% (E)(ii) "the 66,000 measurements" is left out pending your answer on
%%% what the map displays -- see the note at the foot. Restore the count
%%% once it is settled, with the right number.
%%% The single 7.4% became both figures, because the table spans two poses
%%% and the reader needs the one belonging to each row.

\begin{figure}[H]
    \centering

    % First row
    \begin{minipage}[b]{0.45\textwidth}
        \centering
        \includegraphics[width=\textwidth]{Figs/D/wmap_v3_1(nominal).png}

        \vspace{1mm}
        \textbf{(a) Clean scene, iteration 1}
    \end{minipage}
    \hfill
    \begin{minipage}[b]{0.45\textwidth}
        \centering
        \includegraphics[width=\textwidth]{Figs/D/wmap_v3_600(nominal).png}

        \vspace{1mm}
        \textbf{(b) Clean scene, iteration 600}
    \end{minipage}

    \vspace{5mm}

    % Second row
    \begin{minipage}[b]{0.45\textwidth}
        \centering
        \includegraphics[width=\textwidth]{Figs/D/wmap_v3_1(occlusion).png}

        \vspace{1mm}
        \textbf{(c) Occluded scene, iteration 1}
    \end{minipage}
    \hfill
    \begin{minipage}[b]{0.45\textwidth}
        \centering
        \includegraphics[width=\textwidth]{Figs/D/wmap_v3_600(occlusion).png}

        \vspace{1mm}
        \textbf{(d) Occluded scene, iteration 600}
    \end{minipage}

    \caption{Weights $w^{\mathrm{H}}_{j}$ of Variant~3 displayed at their
    image positions, the grey level giving the weight: white a measurement
    retained at $w_{j}=1$, black one rejected at $w_{j}=0$, and the filter
    border excluded. At the first iteration (left column) the clean and
    occluded scenes are almost indistinguishable, the occluder not yet being
    isolated because the residual is large over the whole image and the robust
    scale large in proportion. At convergence (right column) the clean scene
    rejects almost nothing, while the occluded scene rejects the occluder as a
    solid region. The discrimination thus emerges through convergence rather
    than being a precondition for it.}

    \label{fig:wmaps}
\end{figure}

% ======================================================================
%  THREE THINGS TO SETTLE, all quick.
%
%  1. WHAT DOES ONE PIXEL OF THE MAP SHOW? The weights are per
%     measurement, four per position; the map has one value per position.
%     Look at saveMaps() and tell me whether it writes one scale's weight,
%     the mean of the four, or the minimum. Then:
%       - the table caption gets the right count back;
%       - my sentence in 4.6.2 saying these percentages are fractions of P
%         gets corrected if it should be N;
%       - the figure caption can say what is displayed.
%     If it writes the mean of the four, say so explicitly: a reader who
%     assumes it shows a single weight will misread the 0.3% figure.
%
%  2. MEASURE THE BAND. The eight-pixel band around the occluder is twice
%     what the filter support explains. Two possibilities worth
%     distinguishing before you decide what to write:
%       (a) the band is genuinely eight pixels, in which case something
%           other than support is at work and it should be found;
%       (b) the 9.5% and the 5.80% are computed on different bases -- one
%           over positions and the other over measurements, for instance --
%           in which case the ratio 1.6 is an artefact and there may be no
%           excess at all.
%     (b) would be the tidier outcome and is worth ruling out first, since
%     item 1 above may resolve it by itself.
%
%  3. THE G MAPS. My correction to 4.6.4 says the structure measure is
%     recorded "displayed in the same way" and reported in this section.
%     This section shows only weight maps. Either add the G figure here, or
%     tell me it lives in 4.4.2 as fig:structure-maps and I will repoint
%     4.6.4 there. One or the other -- at present 4.6.4 promises a figure
%     that 4.7.3 does not contain.
%
%  --------------------------------------------------------------------
%  CONSEQUENT EDIT TO 4.4.4, arising from item (B)
%
%  The ablation paragraph there argues against the intensity-based measure
%  partly on the ground that "it does not recover the identity when there
%  is nothing to reject: at convergence on a clean target it still rejects
%  1.3% of measurements and holds 8.7% below a quarter weight, against
%  0.3% and 3.2% for the measure adopted". The comparison stands, but the
%  framing does not: we now know the adopted measure does not recover the
%  identity either. Reword to the relative claim, which is the one the
%  numbers support -- that the adopted measure relaxes substantially
%  further as the pose converges, by a factor of four in both statistics,
%  whereas the intensity measure's rejection is largely independent of the
%  pose because brightness does not change as the pose converges. That is
%  a better argument than the one it replaces, and it is now consistent
%  with this section.
%
%  LABELS USED ABOVE
%    sec:nominal            4.7.1
%    sec:occlusion          4.7.2
%    sec:tukey-limitation   4.3.3
%    eq:mad                 the robust scale
%    tab:Occlusion_pose     table of 4.7.2
% ======================================================================
